% !PRINTED 465
\weberpage{465}

We shall call an equation a \gesperrt{normal equation if it is irreducible and has the property that all its roots can be expressed rationally (in $\Omega$) by one of them}. Then a primitive element of a normal field of degree $\mu$ is a root of a normal equation of degree $\mu$. But the converse also holds, namely, that if $\varrho$ is a root of a normal equation, then $\Omega(\varrho)$ is a normal field of the same degree. For if $\varrho_0$ is the root of the normal equation by which all the others are rationally expressed, then $\Omega(\varrho_0)$ is certainly a normal field. But $\Omega(\varrho)$ is conjugate to this field, and hence identical with it. From this it follows further that in a normal equation every root is rationally expressible not only by one determinate root among them, \gesperrt{but by any one whatever}.

One can now reduce arbitrary algebraic fields to normal fields in the following way.

Let
\weq{1}{F(x) = 0}
be an arbitrary reducible or irreducible equation in $\Omega$ of degree $m$, of which we shall only assume \gesperrt{that it has no multiple roots}. Its roots are as follows.
\weq{2}{\alpha,\ \alpha_1 \ldots \alpha_{m-1}.}
From these one obtains $m$ algebraic fields
\weq{3}{\Omega(\alpha),\ \Omega(\alpha_1) \ldots \Omega(\alpha_{m-1}),}
and one can evidently choose the function $F(x)$ so that among the fields (3) any given algebraic fields over $\Omega$ occur. We call the field over $\Omega$ derived from all the quantities (2), i.\,e. from all the roots of the equation (1),
\weq{4}{N = \Omega(\alpha, \alpha_1 \ldots \alpha_{m-1})}
the \gesperrt{Galois field of the equation} $F(x) = 0$.

If $F(x)$ is irreducible, the fields (3) are the fields conjugate to $\Omega(\alpha)$. In this case the field $N$ shall be called the \gesperrt{norm of each of the fields} $\Omega(\alpha), \Omega(\alpha_1) \ldots \Omega(\alpha_{m-1})$.

If $\Omega(\alpha)$ is a normal field, it is identical with its norm. In the general case it is to be shown \gesperrt{that $N$ is a normal field}. We choose a primitive element $\varrho$ of the field $N$ (according to \S.\,143), and can then also denote the field $N$
%CONT>>
